\chapter{Complement Levels}

\section*{What Is This Test?}
The complement system is a group of more than 30 proteins that circulate in the blood and play a critical role in the innate and adaptive immune responses. Complement proteins mediate inflammation, opsonize pathogens for phagocytosis, and directly lyse microbes through the formation of the membrane attack complex (MAC, C5b-9). The complement system is activated through three pathways: the classical pathway (triggered by antigen-antibody complexes and requiring C1q, C1r, C1s, C4, and C2), the lectin pathway (triggered by mannose-binding lectin binding to microbial carbohydrates), and the alternative pathway (triggered by spontaneous C3 hydrolysis on pathogen surfaces). The three pathways converge at C3, leading to formation of the MAC \cite{ricklin2010complement}. Complement testing in clinical practice typically measures C3 and C4 levels, which are the most abundant complement proteins and serve as markers of complement activation and consumption. CH50 (total hemolytic complement) is a functional assay measuring the ability of serum to lyse antibody-coated sheep erythrocytes, requiring all classical pathway components (C1--C9) to be intact. Low complement levels can indicate immune complex-mediated disease (SLE, cryoglobulinemia), congenital complement deficiency, or hepatic synthetic failure.

\section*{Alternative Names}
\begin{itemize}
  \item C3
  \item C4
  \item Complement C3 and C4
  \item CH50 (Total Hemolytic Complement)
  \item AH50 (Alternative Pathway Hemolytic Complement)
  \item Complement Assay
  \item Serum Complement
  \item C3c
  \item C4d
  \item Complement Panel
\end{itemize}
\section*{Principle}
C3 and C4 concentrations are measured by immunoturbidimetry or immunonephelometry on automated chemistry analyzers. Patient serum is incubated with specific anti-C3 or anti-C4 polyclonal antibodies, forming antigen-antibody complexes that scatter light proportionally to complement protein concentration. CH50 is a functional assay: serial dilutions of patient serum are incubated with antibody-coated sheep erythrocytes. If all classical pathway components (C1--C9) are present and functional, the erythrocytes are lysed. The CH50 value is the reciprocal of the serum dilution that produces 50\% hemolysis under standardized conditions. AH50 is a similar functional assay for the alternative pathway, using rabbit erythrocytes (which spontaneously activate the alternative pathway) and blocking the classical and lectin pathways with specific chelators. Individual complement component deficiencies (C1q, C1r, C1s, C2, C4, C3, C5--C9, factor H, factor I) are diagnosed by specialized hemolytic assays or antigenic measurements.

\section*{History}
The complement system was discovered in the 1890s by Jules Bordet, who demonstrated that fresh serum contained a heat-labile factor (``alexine'') that could lyse antibody-coated bacteria \cite{walport2001complement}. Paul Ehrlich coined the term ``complement'' in the 1890s. The classical pathway was the first to be characterized in the early 20th century. The alternative pathway was described by Louis Pillemer in the 1950s. The molecular characterization of complement proteins accelerated in the 1960s--1980s with the purification and sequencing of individual components. The genetic basis of hereditary angioedema (C1 inhibitor deficiency, types I and II) was established in the 1970s--1980s. The association between low complement levels (particularly C4) and active SLE was first recognized in the 1950s. Eculizumab, a monoclonal antibody targeting C5, was approved by the FDA in 2007 for paroxysmal nocturnal hemoglobinuria (PNH) and later for atypical hemolytic uremic syndrome (aHUS), myasthenia gravis, and neuromyelitis optica spectrum disorder, ushering in the era of therapeutic complement inhibition \cite{walport2001complement2}.

\section*{Why This Test Is Ordered}
\begin{itemize}
  \item Evaluation and monitoring of systemic lupus erythematosus (SLE) activity --- low C3 and C4 indicate active immune complex disease with complement consumption, particularly active lupus nephritis \cite{truedsson2007complement}
  \item Diagnosis and monitoring of mixed cryoglobulinemia (associated with chronic hepatitis C virus) with low C4 and normal C3
  \item Investigation of recurrent infections with encapsulated organisms (Streptococcus pneumoniae, Neisseria meningitidis, Haemophilus influenzae) suggestive of congenital complement component deficiency (particularly C2, C3, C5--C9)
  \item Diagnosis of hereditary angioedema (HAE) --- C4 is low during and between attacks due to uncontrolled C1 activation; C1 inhibitor level and function must be measured for confirmation
  \item Evaluation of membranoproliferative glomerulonephritis (MPGN), particularly type II (dense deposit disease) and C3 glomerulopathy, which are driven by alternative pathway dysregulation
  \item Investigation of angioedema without urticaria --- low C4 and C1 inhibitor deficiency are characteristic of hereditary angioedema, while angiotensin-converting enzyme (ACE) inhibitor-induced angioedema has normal complement
  \item Monitoring of disease activity and treatment response in SLE, post-streptococcal glomerulonephritis, and cryoglobulinemia
  \item Evaluation of atypical hemolytic uremic syndrome (aHUS) --- excessive alternative pathway activation leads to thrombotic microangiopathy with low C3 and normal C4
\end{itemize}

\section*{Primary vs Secondary Test}
This is a secondary test used in conjunction with specific disease markers. In SLE, C3 and C4 are used primarily for monitoring disease activity and are part of the SLEDAI (SLE Disease Activity Index) and other disease activity scoring systems. For complement deficiency syndromes, complement measurement is a primary diagnostic test.

\section*{How This Test Is Performed}
For most patients, a health care professional will take a blood sample from a vein in your arm using a small needle. After the needle is inserted, a small amount of blood will be collected into a test tube or vial. You may feel a little sting when the needle goes in or out. This usually takes less than five minutes.

\section*{What Preparation Is Needed}
No special preparation is required. Fasting is not necessary. For CH50 and AH50 functional assays, the specimen should be processed promptly (within 2 hours of collection) and the serum should be frozen at -70 degrees C if not tested immediately, as complement components are heat-labile and degrade rapidly at room temperature. Collection in a serum separator tube (SST, gold-top) or red-top tube is required; EDTA plasma interferes with functional hemolytic assays by chelating calcium and magnesium required for complement activation.

\section*{Contraindications and Precautions}
\begin{itemize}
  \item No absolute contraindications. Improper specimen handling is the most common cause of spurious complement results. Complement proteins are heat-labile --- delays in processing or prolonged storage at room temperature cause falsely low results. Inappropriate specimen type (EDTA plasma rather than serum) yields falsely low CH50 due to calcium and magnesium chelation. C3 and C4 are acute phase reactants produced by the liver, and their levels can be influenced by hepatic synthetic function --- end-stage liver disease reduces C3 and C4 production, and acute inflammation may increase their production. C3 and C4 are typically measured simultaneously because C4 is more sensitive to immune complex consumption (classical pathway activation) while C3 reflects both classical and alternative pathway activation. Genetic complement deficiencies (C1, C2, C3, C4, C5--C9) should be diagnosed during quiescent disease, not during active inflammation which may consume and lower complement levels.
\end{itemize}
\section*{Risks}
Minimal risk. Slight pain or bruising at the venipuncture site. Rare: hematoma, infection, vasovagal syncope.

\section*{What Happens After the Test}
Results are typically available within 1 to 2 hours for C3 and C4, and within 1 to 3 days for CH50. Low C3 and C4 in a patient with known SLE suggest active disease and may predict a lupus nephritis flare. Persistently low C4 with normal C3 raises suspicion for mixed cryoglobulinemia. Low CH50 in a patient with recurrent Neisseria infections suggests a terminal complement component deficiency (C5--C9). Low C4 with normal C3 and a history of recurrent angioedema warrants C1 inhibitor testing for hereditary angioedema. Your healthcare provider will interpret complement results in the context of clinical presentation and other laboratory findings.

\section*{Understanding Results}

\subsection*{Below Normal}
\begin{itemize}
  \item Active systemic lupus erythematosus with immune complex-mediated complement consumption: Low C3 AND low C4 (classical pathway activation). Low C3 with low or normal C4 suggests alternative pathway activation, which can occur with C3 nephritic factor
  \item Active lupus nephritis: Markedly low C3 and C4 indicate active immune complex glomerulonephritis. Rising complement levels indicate clinical response to treatment. C4 levels recover more rapidly than C3
  \item Cryoglobulinemia (particularly type II mixed cryoglobulinemia with hepatitis C virus): Persistently low C4 with relatively normal C3 is characteristic, as cryoglobulins preferentially activate the classical pathway leading to C4 consumption
  \item Post-streptococcal (post-infectious) glomerulonephritis: Low C3 with normal C4 (alternative pathway activation) lasting 6--8 weeks. Normalization of C3 is expected with clinical recovery
  \item Membranoproliferative glomerulonephritis (MPGN) type II/dense deposit disease and C3 glomerulopathy: Persistent low C3 with normal C4 due to autoantibody (C3 nephritic factor, C3NeF) that stabilizes the alternative pathway C3 convertase, leading to continuous C3 consumption \cite{pickering2002uncontrolled}
  \item Hereditary angioedema (HAE): Low C4 between and during attacks, low C1 inhibitor antigen and/or function. C3 is typically normal
  \item Congenital complement component deficiency: Homozygous deficiency of classical pathway components (C1q, C1r, C1s, C4, C2) is associated with autoimmune disease (particularly SLE-like syndrome). C5--C9 and properdin deficiency predisposes to recurrent Neisseria infections. C3 deficiency (very rare) leads to recurrent severe pyogenic infections
  \item End-stage liver disease: Impaired hepatic synthesis of complement proteins leads to low C3 and C4. The low complement reflects impaired protein synthesis rather than immune consumption
  \item Severe protein malnutrition or protein-losing enteropathy/nephropathy: Loss of complement proteins results in low C3 and C4
  \item Eculizumab (anti-C5) therapy: Pharmacologic C5 blockade prevents C5--C9 assembly and MAC formation. CH50 is zero or extremely low; C3 and C4 levels are typically normal
\end{itemize}

\subsection*{Above Normal}
\begin{itemize}
  \item Acute phase response: C3 and C4 are acute phase reactants. Elevated C3 and C4 are seen in acute infection, inflammation, trauma, surgery, and malignancy
  \item Active systemic inflammation: C3 and C4 can be elevated in acute infections (particularly bacterial), inflammatory conditions (rheumatoid arthritis, inflammatory bowel disease), and after myocardial infarction
  \item Pregnancy: C3 and C4 levels increase by 20--50\% during the third trimester due to increased hepatic synthesis stimulated by estrogen
  \item Obstructive jaundice: Biliary obstruction leads to elevated complement protein synthesis by the liver
  \item Elevated C3 with low CH50 may indicate an acute phase response with a concurrent complement deficiency in another component --- CH50 reflects overall classical pathway function and should be measured if isolated elevation is found
\end{itemize}

\section*{Normal Results}
\begin{itemize}
  \item \textbf{C3 (Complement 3):} 80--180 mg/dL (adults, age-dependent reference ranges apply for children)
  \item \textbf{C4 (Complement 4):} 15--45 mg/dL (adults, age-dependent reference ranges apply for children)
  \item \textbf{CH50 (Total Hemolytic Complement):} 60--144 CH50 Units/mL or 30--75\% of normal serum complement activity
  \item \textbf{AH50 (Alternative Pathway Hemolytic Complement):} >46 US units/mL or >75\% of reference serum activity
  \item \textbf{Pregnancy (third trimester):} C3 and C4 increase by 20--50\% above non-pregnant levels
  \item \textbf{Newborns and infants:} Complement levels are lower at birth (C3 approximately 50\%, C4 approximately 60\% of adult levels) and reach adult levels by 6--12 months of age
  \item Interpretation requires simultaneous consideration of C3, C4, and clinical context, ideally with a CH50 functional assay when complement deficiency is suspected
\end{itemize}

\section*{Follow-up Tests}
\begin{itemize}
  \item \textbf{CH50 (Total Hemolytic Complement):} Functional assay of the complete classical complement pathway (C1--C9). A CH50 of zero in a patient with normal C3 and C4 antigen levels suggests a deficiency in an individual complement component that should be further investigated
  \item \textbf{AH50 (Alternative Pathway Hemolytic Complement):} Functional assay of the alternative pathway. Low AH50 with normal CH50 suggests alternative pathway dysfunction (properdin deficiency, factor H or factor I deficiency, C3 nephritic factor)
  \item \textbf{C1 inhibitor (C1-INH) antigen and function:} If hereditary angioedema is suspected (recurrent angioedema without urticaria, low C4, positive family history). Type I HAE has low C1-INH antigen and function; type II HAE has normal or elevated C1-INH antigen but low function
  \item \textbf{C3 nephritic factor (C3NeF):} Autoantibody that stabilizes the alternative pathway C3 convertase, causing persistent C3 consumption. Tested if persistent low C3 with normal C4 suggests C3 glomerulopathy or MPGN
  \item \textbf{Individual complement component levels (C1q, C2, C5, C6, C7, C8, C9):} If total complement deficiency syndrome is suspected (recurrent neisserial infections, autoimmune disease)
  \item \textbf{Antinuclear antibody (ANA) and anti-dsDNA antibody:} If low C3 and C4 suggest active SLE. Rising anti-dsDNA titers with falling C3 and C4 predict lupus nephritis flares
  \item \textbf{Rheumatoid factor, cryoglobulins, hepatitis C serology:} If persistently low C4 with normal C3 suggests cryoglobulinemia
  \item \textbf{Anti-streptolysin O (ASO) and anti-DNase B titers:} If low C3 with normal C4 suggests post-streptococcal glomerulonephritis
  \item \textbf{Urinalysis and spot urine protein-to-creatinine ratio:} If glomerulonephritis is suspected, to detect hematuria, red blood cell casts, and proteinuria
  \item \textbf{Serum creatinine and eGFR:} Assessment of renal function, particularly if lupus nephritis, cryoglobulinemia-associated nephritis, or MPGN is suspected
\end{itemize}

\section*{References}
\bibliographystyle{unsrtnat}
\bibliography{references}

\clearpage
\section*{Regulatory Status and Authorization Date}
This chapter describes a test category, not a specific commercial product. FDA, CDSCO, EU IVDR, and MHRA approval or registration dates therefore cannot be assigned without the assay manufacturer, product name, instrument, and intended use. For laboratory-developed tests, an individual agency approval date may not exist. See the Preface for the interpretation of regulatory status.

\clearpage
